% GENERATED FILE. Edit canonical lessons, then run npm run generate:books.
% canonical-source-hash: fnv1a64:4d2976e4da9baeee
% canonical-lessons: PA-C333-nai, PA-C333-munnna, PA-C333-dhobi, PA-C333-mochi, PA-C333-mali

\chapter{The Barber and the Washerman}
\label{ch:pa-a2-the-barber-and-the-washerman}
\hlchaptermodality{333}

\begin{chapteropening}
\textbf{By the end of this chapter:} \emph{I can name the people who cut hair and wash clothes.}

Complete the last lesson of chapter 333: I can name the people who cut hair and wash clothes.
\end{chapteropening}

\section[\texorpdfstring{\pa{ਨਾਈ} (nāī)}{ਨਾਈ (nāī)}]{\pa{ਨਾਈ} (nāī) — a barber}

\label{lesson:PA-C333-nai}



\begin{quote}
Before the new one: say the Punjabi for a woman's female friend, then the Punjabi for to call out to someone.
\end{quote}

\subsection*{You'll want to know: \pa{ਨਾਈ}}
\textbf{\pa{ਨਾਈ}} — \emph{nāī} — \textquotedblleft{}a barber\textquotedblright{}.

\begin{quote}
\pa{ਨਾਈ} \pa{ਹਰ} \pa{ਸਵੇਰ} \pa{ਆਪਣੀ} \pa{ਦੁਕਾਨ} \pa{ਖੋਲ੍ਹਦਾ} \pa{ਹੈ।}
\end{quote}

\begin{quote}
\emph{nāī har saver āpṇī dukān kholhdā hai.}
\end{quote}

\begin{quote}
The barber opens his shop every morning.
\end{quote}

A masculine noun.

\begin{cousinweb}
\textbf{\pa{ਨਾਈ}} (\emph{nāī}) is inherited from Sanskrit \emph{nāpita}, \textquotedblleft{}a barber\textquotedblright{}; CDIAL lists Punjabi \emph{nāī}, \textquotedblleft{}barber\textquotedblright{}, under it (CDIAL 13792).
\end{cousinweb}

\begin{grammarlens}[title={What you've built}]
One more word for this chapter.
\end{grammarlens}

\subsection*{Your turn}
\begin{itemize}
\raggedright
  \item \emph{Say it:} \emph{nāī}
  \item \emph{Say it:} \emph{nāī}, once more
  \item \emph{Say it:} \emph{nāī har saver āpṇī dukān kholhdā hai.} — The barber opens his shop every morning.
  \item \emph{Recall:} say the Punjabi for a woman's female friend, then the Punjabi for to call out to someone, then say \emph{nāī} again
\end{itemize}

\begin{tcolorbox}[breakable,colback=teal!4,colframe=teal!35!black,title={Before you move on}]
What does \pa{ਨਾਈ} mean? (\textquotedblleft{}A barber\textquotedblright{}.) Say it once more.
\end{tcolorbox}
\section[\texorpdfstring{\pa{ਮੁੰਨਣਾ} (munnṇā)}{ਮੁੰਨਣਾ (munnṇā)}]{\pa{ਮੁੰਨਣਾ} (munnṇā) — to shave}

\label{lesson:PA-C333-munnna}



\begin{quote}
Before the new one: say the Punjabi for to call out to someone, then the Punjabi for a barber.
\end{quote}

\subsection*{You'll want to know: \pa{ਮੁੰਨਣਾ}}
\textbf{\pa{ਮੁੰਨਣਾ}} — \emph{munnṇā} — \textquotedblleft{}to shave\textquotedblright{}.

\begin{quote}
\pa{ਨਾਈ} \pa{ਸਵੇਰੇ} \pa{ਲੋਕਾਂ} \pa{ਦੀ} \pa{ਦਾੜ੍ਹੀ} \pa{ਮੁੰਨਦਾ} \pa{ਹੈ।}
\end{quote}

\begin{quote}
\emph{nāī savere lokāṁ dī dāṛhī munndā hai.}
\end{quote}

\begin{quote}
The barber shaves people's beards in the morning.
\end{quote}

Past (I …): \pa{ਮੈਂ} \pa{ਆਪਣੀ} \pa{ਦਾੜ੍ਹੀ} \pa{ਮੁੰਨੀ}, \emph{maiṁ āpṇī dāṛhī munnī}.

\begin{cousinweb}
\textbf{\pa{ਮੁੰਨਣਾ}} (\emph{munnṇā}) is inherited from Sanskrit \emph{muṇḍayati}, \textquotedblleft{}shaves\textquotedblright{}; CDIAL lists Punjabi \emph{munnanā}, \textquotedblleft{}shaves\textquotedblright{}, under it (CDIAL 10194).
\end{cousinweb}

\begin{grammarlens}[title={What you've built}]
One more word for this chapter.
\end{grammarlens}

\subsection*{Your turn}
\begin{itemize}
\raggedright
  \item \emph{Say it:} \emph{munnṇā}
  \item \emph{Say it:} \emph{munnṇā}, once more
  \item \emph{Say it:} \emph{nāī savere lokāṁ dī dāṛhī munndā hai.} — The barber shaves people's beards in the morning.
  \item \emph{Recall:} say the Punjabi for to call out to someone, then the Punjabi for a barber, then say \emph{munnṇā} again
  \item \emph{Write it:} copy \textbf{\pa{ਮੁੰਨਣਾ}} once, checking the whole word
\end{itemize}

\begin{tcolorbox}[breakable,colback=teal!4,colframe=teal!35!black,title={Before you move on}]
What does \pa{ਮੁੰਨਣਾ} mean? (\textquotedblleft{}To shave\textquotedblright{}.) Say it once more.
\end{tcolorbox}
\section[\texorpdfstring{\pa{ਧੋਬੀ} (dhobī)}{ਧੋਬੀ (dhobī)}]{\pa{ਧੋਬੀ} (dhobī) — a washerman}

\label{lesson:PA-C333-dhobi}



\begin{quote}
Before the new one: say the Punjabi for a barber, then the Punjabi for to shave.
\end{quote}

\subsection*{You'll want to know: \pa{ਧੋਬੀ}}
\textbf{\pa{ਧੋਬੀ}} — \emph{dhobī} — \textquotedblleft{}a washerman\textquotedblright{}.

\begin{quote}
\pa{ਧੋਬੀ} \pa{ਹਰ} \pa{ਹਫ਼ਤੇ} \pa{ਸਾਡੇ} \pa{ਕੱਪੜੇ} \pa{ਧੋਂਦਾ} \pa{ਹੈ।}
\end{quote}

\begin{quote}
\emph{dhobī har hafte sāḍe kappṛe dhondā hai.}
\end{quote}

\begin{quote}
The washerman washes our clothes every week.
\end{quote}

A masculine noun.

\begin{cousinweb}
\textbf{\pa{ਧੋਬੀ}} (\emph{dhobī}) goes back to a reconstructed \emph{*dhauvati}, \textquotedblleft{}washes\textquotedblright{}, as a word taken from Hindi: CDIAL lists Punjabi \emph{dhobī} under it among the forms that spread from Hindi \emph{dhobī}, \textquotedblleft{}washerman\textquotedblright{} (CDIAL 6886).
\end{cousinweb}

\begin{grammarlens}[title={What you've built}]
One more word for this chapter.
\end{grammarlens}

\subsection*{Your turn}
\begin{itemize}
\raggedright
  \item \emph{Say it:} \emph{dhobī}
  \item \emph{Say it:} \emph{dhobī}, once more
  \item \emph{Say it:} \emph{dhobī har hafte sāḍe kappṛe dhondā hai.} — The washerman washes our clothes every week.
  \item \emph{Recall:} say the Punjabi for a barber, then the Punjabi for to shave, then say \emph{dhobī} again
  \item \emph{Write it:} copy \textbf{\pa{ਧੋਬੀ}} once, checking the whole word
\end{itemize}

\begin{tcolorbox}[breakable,colback=teal!4,colframe=teal!35!black,title={Before you move on}]
What does \pa{ਧੋਬੀ} mean? (\textquotedblleft{}A washerman\textquotedblright{}.) Say it once more.
\end{tcolorbox}
\section[\texorpdfstring{\pa{ਮੋਚੀ} (mochī)}{ਮੋਚੀ (mochī)}]{\pa{ਮੋਚੀ} (mochī) — a cobbler}

\label{lesson:PA-C333-mochi}



\begin{quote}
Before the new one: say the Punjabi for to shave, then the Punjabi for a washerman.
\end{quote}

\subsection*{You'll want to know: \pa{ਮੋਚੀ}}
\textbf{\pa{ਮੋਚੀ}} — \emph{mochī} — \textquotedblleft{}a cobbler\textquotedblright{}.

\begin{quote}
\pa{ਮੋਚੀ} \pa{ਬਾਜ਼ਾਰ} \pa{ਵਿੱਚ} \pa{ਜੁੱਤੀਆਂ} \pa{ਸਿਉਂਦਾ} \pa{ਹੈ।}
\end{quote}

\begin{quote}
\emph{mochī bāzār vich juttīāṁ siundā hai.}
\end{quote}

\begin{quote}
The cobbler stitches shoes in the market.
\end{quote}

A masculine noun.

\begin{cousinweb}
\textbf{\pa{ਮੋਚੀ}} (\emph{mochī}) is inherited from Sanskrit \emph{mocika}, \textquotedblleft{}a tanner and cobbler\textquotedblright{}; CDIAL lists Punjabi \emph{mocī}, \textquotedblleft{}cobbler\textquotedblright{}, under it (CDIAL 10349). CDIAL traces the word for a shoe behind it to Iranian.
\end{cousinweb}

\begin{grammarlens}[title={What you've built}]
One more word for this chapter.
\end{grammarlens}

\subsection*{Your turn}
\begin{itemize}
\raggedright
  \item \emph{Say it:} \emph{mochī}
  \item \emph{Say it:} \emph{mochī}, once more
  \item \emph{Say it:} \emph{mochī bāzār vich juttīāṁ siundā hai.} — The cobbler stitches shoes in the market.
  \item \emph{Recall:} say the Punjabi for to shave, then the Punjabi for a washerman, then say \emph{mochī} again
  \item \emph{Write it:} copy \textbf{\pa{ਮੋਚੀ}} once, checking the whole word
\end{itemize}

\begin{tcolorbox}[breakable,colback=teal!4,colframe=teal!35!black,title={Before you move on}]
What does \pa{ਮੋਚੀ} mean? (\textquotedblleft{}A cobbler\textquotedblright{}.) Say it once more.
\end{tcolorbox}
\section[\texorpdfstring{\pa{ਮਾਲੀ} (mālī)}{ਮਾਲੀ (mālī)}]{\pa{ਮਾਲੀ} (mālī) — a gardener}

\label{lesson:PA-C333-mali}



\begin{quote}
Before the new one: say the Punjabi for a washerman, then the Punjabi for a cobbler.
\end{quote}

\subsection*{You'll want to know: \pa{ਮਾਲੀ}}
\textbf{\pa{ਮਾਲੀ}} — \emph{mālī} — \textquotedblleft{}a gardener\textquotedblright{}.

\begin{quote}
\pa{ਮਾਲੀ} \pa{ਸਵੇਰੇ} \pa{ਬਾਗ਼} \pa{ਵਿੱਚ} \pa{ਪਾਣੀ} \pa{ਪਾਉਂਦਾ} \pa{ਹੈ।}
\end{quote}

\begin{quote}
\emph{mālī savere bāgh vich pāṇī pāundā hai.}
\end{quote}

\begin{quote}
The gardener waters the garden in the morning.
\end{quote}

A masculine noun.

\begin{cousinweb}
\textbf{\pa{ਮਾਲੀ}} (\emph{mālī}) is inherited from Sanskrit \emph{mālin}, \textquotedblleft{}garlanded; a gardener\textquotedblright{}; CDIAL lists Punjabi \emph{mālī}, \textquotedblleft{}garlanded; gardener\textquotedblright{}, under it (CDIAL 10094).
\end{cousinweb}

\begin{grammarlens}[title={What you've built}]
One more word for this chapter.
\end{grammarlens}

\subsection*{Your turn}
\begin{itemize}
\raggedright
  \item \emph{Say it:} \emph{mālī}
  \item \emph{Say it:} \emph{mālī}, once more
  \item \emph{Say it:} \emph{mālī savere bāgh vich pāṇī pāundā hai.} — The gardener waters the garden in the morning.
  \item \emph{Recall:} say the Punjabi for a washerman, then the Punjabi for a cobbler, then say \emph{mālī} again
  \item \emph{Write it:} copy \textbf{\pa{ਮਾਲੀ}} once, checking the whole word
\end{itemize}

\begin{tcolorbox}[breakable,colback=teal!4,colframe=teal!35!black,title={Before you move on}]
What does \pa{ਮਾਲੀ} mean? (\textquotedblleft{}A gardener\textquotedblright{}.) Say it once more.
\end{tcolorbox}
